\documentclass[12pt,a4paper]{article}
\usepackage[UTF8]{ctex}
\usepackage{amsmath,amssymb}
\usepackage{geometry}
\usepackage{xcolor}
\usepackage{listings}
\usepackage{tcolorbox}
\geometry{margin=2cm}

\lstset{
  language=C,
  basicstyle=\ttfamily\small,
  keywordstyle=\color{blue}\bfseries,
  commentstyle=\color{gray},
  numbers=left,
  numberstyle=\tiny\color{gray},
  frame=single,
  breaklines=true,
  showstringspaces=false,
  tabsize=2
}

\newtcolorbox{bugbox}[1][]{colback=red!5!white,colframe=red!75!black,title=常见错误,#1}
\newtcolorbox{keybox}[1][]{colback=yellow!5!white,colframe=yellow!75!black,title=考点,#1}

\title{FDS 复习 — 排序算法}
\author{6月29日 冲刺}
\date{}

\begin{document}
\maketitle

\section{排序算法总览}

\begin{tabular}{|c|c|c|c|c|c|}
\hline
\textbf{算法} & \textbf{平均} & \textbf{最坏} & \textbf{空间} & \textbf{稳定} & \textbf{比较} \\
\hline
插入排序 & $O(N^2)$ & $O(N^2)$ & $O(1)$ & 稳定 & 是 \\
希尔排序 & $O(N^{3/2})$ & $O(N^2)$ & $O(1)$ & 不稳 & 是 \\
堆排序 & $O(N\log N)$ & $O(N\log N)$ & $O(1)$ & 不稳 & 是 \\
归并排序 & $O(N\log N)$ & $O(N\log N)$ & $O(N)$ & 稳定 & 是 \\
快速排序 & $O(N\log N)$ & $O(N^2)$ & $O(\log N)$ & 不稳 & 是 \\
桶排序 & $O(M+N)$ & — & $O(M)$ & 稳定 & 否 \\
基数排序 & $O(P(N+B))$ & — & $O(N+B)$ & 稳定 & 否 \\
\hline
\end{tabular}

\begin{keybox}
\textbf{比较排序下界：} 任意基于比较的排序算法，最坏情况下至少需要 $\Omega(N\log N)$ 次比较。

桶排和基数排不基于比较，因此能突破 $O(N\log N)$ 下界。
\end{keybox}

\section{插入排序 (Insertion Sort)}

\textbf{思想：} 像整理扑克牌，每次从后面取一张插入前面已排好的位置。

\begin{lstlisting}
void insertionSort(int A[], int N) {
    int i, j, tmp;
    for (i = 1; i < N; i++) {
        tmp = A[i];
        for (j = i; j > 0 && A[j-1] > tmp; j--)
            A[j] = A[j-1];        // 比tmp大的右移
        A[j] = tmp;               // 插入到正确位置
    }
}
\end{lstlisting}

\textbf{复杂度：}
\begin{itemize}
\item 最好（已有序）：$O(N)$ — 内层循环一次不进
\item 最坏（逆序）：$O(N^2)$
\item 平均逆序数：$N(N-1)/4$
\end{itemize}

\begin{bugbox}
\begin{itemize}
\item 第一次写成了 \texttt{tmp = A[i]} 后直接从头找位置，导致不比较大小
\item 交换写成 \texttt{A[j]=A[j-1]} 忘了 \texttt{A[j-1]=A[j]} ——应该只是右移不是交换
\end{itemize}
\end{bugbox}

\section{希尔排序 (Shellsort)}

\textbf{思想：} 不等距的插入排序。先用大间隔跳着排，逐步缩小间隔，最后间隔1 = 普通插入排序。

\begin{lstlisting}
void shellSort(int A[], int N) {
    int gap, i, j, tmp;
    for (gap = N/2; gap > 0; gap /= 2) {
        for (i = gap; i < N; i++) {
            tmp = A[i];
            for (j = i; j >= gap && A[j-gap] > tmp; j -= gap)
                A[j] = A[j-gap];
            A[j] = tmp;
        }
    }
}
\end{lstlisting}

\textbf{和插入排序的区别（只改3处）：}
\begin{enumerate}
\item 外套一层 \texttt{for(gap=N/2; gap>0; gap/=2)}
\item \texttt{j-1} 变成 \texttt{j-gap}
\item \texttt{j--} 变成 \texttt{j-=gap}
\end{enumerate}

\begin{bugbox}
\begin{itemize}
\item \texttt{while(gap!=1)} 吞了 gap=1 这轮，最后一轮常规插入没跑
\item \texttt{i += gap} 只处理了一组，其他组的全漏了
\item 内层 for 里声明的 j 在外面 \texttt{A[j]=tmp} 用不了
\end{itemize}
\end{bugbox}

\textbf{Hibbard 序列：} $h_k = 2^k - 1$，最坏 $\Theta(N^{3/2})$。

\section{堆排序 (Heapsort)}

\textbf{思想：} 原地最大堆 + 反复取堆顶放到末尾。

\begin{itemize}
\item 索引从 0 开始！ left = 2i+1, right = 2i+2, parent = (i-1)/2
\item 用\textbf{最大堆}：根是最大的，每次把根和末尾交换后下滤
\end{itemize}

\begin{lstlisting}
void percDown(int A[], int i, int N) {
    int child, tmp = A[i];
    while (2*i + 1 < N) {            // 有左孩子
        child = 2*i + 1;
        if (child+1 < N && A[child+1] > A[child])  // 选较大孩子
            child++;
        if (tmp >= A[child]) break;   // 满足堆序
        A[i] = A[child];              // 孩子上移
        i = child;
    }
    A[i] = tmp;                       // 放到最终位置
}

void heapSort(int A[], int N) {
    int i;
    for (i = N/2 - 1; i >= 0; i--)    // BuildMaxHeap: O(N)
        percDown(A, i, N);
    for (i = N - 1; i > 0; i--) {     // 排序: O(N log N)
        int tmp = A[0]; A[0] = A[i]; A[i] = tmp;
        percDown(A, 0, i);
    }
}
\end{lstlisting}

\begin{bugbox}
\begin{itemize}
\item 忘了堆从 0 开始，孩子/父公式写成 \texttt{2*i} 和 \texttt{i/2}
\item \texttt{j > 1} 漏了根，应该是 \texttt{j >= 1}（ADT版）或 \texttt{j >= 0}（原地版）
\item \texttt{h->size} 和 \texttt{A[child]} 语法混用
\end{itemize}
\end{bugbox}

\textbf{复杂度：} BuildHeap $O(N)$ + $N$次DeleteMax $O(N\log N)$ = $O(N\log N)$

\section{归并排序 (Mergesort)}

\textbf{思想：} 分治法。拆成两半递归排，再合并两个有序数组。

\begin{lstlisting}
void merge(int A[], int left, int mid, int right) {
    int i = left, j = mid + 1, k = 0;
    int n = right - left + 1;
    int* tmp = (int*)malloc(n * sizeof(int));

    while (i <= mid && j <= right) {
        if (A[i] <= A[j]) tmp[k++] = A[i++];
        else              tmp[k++] = A[j++];
    }
    while (i <= mid)   tmp[k++] = A[i++];
    while (j <= right) tmp[k++] = A[j++];
    for (int p = 0; p < n; p++)
        A[left + p] = tmp[p];
    free(tmp);
}

void mergeSort(int A[], int left, int right) {
    if (left >= right) return;
    int mid = (left + right) / 2;
    mergeSort(A, left, mid);
    mergeSort(A, mid + 1, right);
    merge(A, left, mid, right);
}
\end{lstlisting}

\begin{bugbox}
\begin{itemize}
\item tmp 的 k 从 left 开始（应该是从0），拷回时搞混
\item 忘了 free(tmp)
\item n 声明在 malloc 后面，编译错误
\end{itemize}
\end{bugbox}

\section{快速排序 (Quicksort)}

\textbf{思想：} 选 pivot，小的放左，大的放右，递归。

\begin{lstlisting}
int partition(int A[], int left, int right) {
    int pivot = A[right];         // 选最右边为pivot
    int i = left;                 // i：小元素边界

    for (int j = left; j < right; j++) {
        if (A[j] < pivot) {
            int tmp = A[i]; A[i] = A[j]; A[j] = tmp;
            i++;
        }
    }
    int tmp = A[i]; A[i] = A[right]; A[right] = tmp;
    return i;                     // pivot最终位置
}

void quickSort(int A[], int left, int right) {
    if (left >= right) return;
    int p = partition(A, left, right);
    quickSort(A, left, p - 1);
    quickSort(A, p + 1, right);
}
\end{lstlisting}

\textbf{手算示例：} A = [5,2,6,1,4,9,3,7], pivot=7

\begin{verbatim}
j=0: 5<7 → swap 5↔5  i=1  [5,2,6,1,4,9,3,7]
j=1: 2<7 → swap 2↔2  i=2  [5,2,6,1,4,9,3,7]
j=2: 6<7 → swap 6↔6  i=3  [5,2,6,1,4,9,3,7]
j=3: 1<7 → swap 1↔1  i=4  [5,2,6,1,4,9,3,7]
j=4: 4<7 → swap 4↔4  i=5  [5,2,6,1,4,9,3,7]
j=5: 9<7? NO           i=5  [5,2,6,1,4,9,3,7]
j=6: 3<7 → swap 9↔3  i=6  [5,2,6,1,4,3,9,7]
最后 swap(7,9) → [5,2,6,1,4,3,7,9]
pivot=7在位置6，左全<7，右全>7
\end{verbatim}

\textbf{优化：} Median-of-Three — 取 left/mid/right 三元素的中值作为 pivot。

\begin{keybox}
\textbf{最坏情况：} 每次 pivot 选到最小或最大 → $O(N^2)$\\
\textbf{平均：} $O(N\log N)$
\end{keybox}

\section{基数排序 (Radix Sort — LSD)}

\textbf{思想：} 不比较，按每一位丢桶里排序。从最低位(LSD)到最高位(MSD)。

\textbf{示例：} [64,8,216,512,27,729,0,1,343,125]

\begin{verbatim}
Pass1(个位): 0:[0] 1:[1] 2:[512] 3:[343] 4:[64] 5:[125] 6:[216] 7:[27] 8:[8] 9:[729]
→ 0,1,512,343,64,125,216,27,8,729

Pass2(十位): 0:[0,1,8] 1:[512,216] 2:[125,27,729] 3:[343] 6:[64]
→ 0,1,8,512,216,125,27,729,343,64

Pass3(百位): 0:[0,1,8,27,64] 1:[125] 2:[216] 3:[343] 5:[512] 7:[729]
→ 0,1,8,27,64,125,216,343,512,729 ✓
\end{verbatim}

\textbf{复杂度：} $T = O(P \cdot (N + B))$，P=位数，B=基数。

\section{表排序 (Table Sort)}

\textbf{动机：} 排序大结构体时，移动代价高。只排序索引表。

\begin{verbatim}
原始: A = [f, d, a, b, e, c]
T:    [0, 1, 2, 3, 4, 5]
排序后: T = [2, 3, 5, 1, 4, 0]
  T[0]=2 → A[2]=a (最小)
\end{verbatim}

\textbf{物理重排：} 利用置换环，最坏 $\lfloor 3N/2 \rfloor$ 次移动。

\begin{keybox}
\textbf{考点：}
\begin{itemize}
\item 稳定 vs 不稳定（考试必问！）
\item 各算法的最坏/平均/空间复杂度
\item 为什么桶排/基数排能突破 $O(N\log N)$ —— 不基于比较
\item 决策树 → 比较排序下界 $\Omega(N\log N)$
\item 快排 pivot 选不好 → $O(N^2)$
\end{itemize}
\end{keybox}

\end{document}
